%%% Кодировки и шрифты %%%
% Язык по-умолчанию русский с поддержкой приятных команд пакета babel
\setmainlanguage[babelshorthands=true]{russian}
% Дополнительный язык = английский (в американской вариации по-умолчанию)
\setotherlanguage{english}

\setmonofont{Source Code Pro}[Scale=0.9]                    % Моноширинный шрифт
\newfontfamily\cyrillicfonttt{Source Code Pro}[Scale=0.9]   % Моноширинный шрифт для кириллицы

\defaultfontfeatures{Ligatures=TeX}                 % Стандартные лигатуры TeX, замены нескольких дефисов на тире и т. п. Настройки моноширинного шрифта должны идти до этой строки, чтобы при врезках кода программ в коде не применялись лигатуры и замены дефисов.

\setmainfont{STIX Two Text}                         % Шрифт с засечками
\newfontfamily\cyrillicfont{STIX Two Text}          % Шрифт с засечками для кириллицы

\setsansfont{Source Sans Pro}                       % Шрифт без засечек
\newfontfamily\cyrillicfontsf{Source Sans Pro}      % Шрифт без засечек для кириллицы

\setmathfont{STIX Two Math}[StylisticSet=8]         % Математический шрифт
\newfontfamily\cyrillicfontmf{STIX Two Math}        % Математический шрифт для кириллицы
